% Cote 119: offprints, reports, courses, CV and letters (typed and printed part): archive cote 119, site page 70 (PDF page 71).
% offprint. Draft compared: none.
% AI-transcribed by Claude Opus 5.5 (claude-opus-5-5, adaptive thinking, effort high), Message Batches API, batch r119_full (msgbatch_01DwZNm5Y7YsKBWyAN4sFs8w), 2026-11-13; not yet checked by eye.
\scanpage{119/70}
\mspage{( 2 )}

qui associe à $g\in \mathrm{G}(\mathrm{S}')$ le $\otimes$-endomorphisme de $(\omega^{\mathrm{G}})_{\mathrm{S}'}$ induisant sur chaque
G-$\mathcal{O}_{\mathrm{S}}$-Module $\mathcal{F}$ la multiplication par $g$ dans $\omega^{\mathrm{G}}(\mathcal{F})\otimes_{\mathcal{O}_{\mathrm{S}}}\mathcal{O}_{\mathrm{S}'}$.

Proposition 1.3. --- Soient S un schéma, G un S-monoïde affine. Alors,
(1.2.1) est un isomorphisme.

On en déduit aisément :

Corollaire 1.4. --- Soient G, G$'$ des S-monoïdes affines. L'application
$u\mapsto \mathrm{Rep}(u)$ induit une bijection de l'ensemble $\mathrm{Hom}_{\text{S-mon}}(\mathrm{G},\mathrm{G}')$ des
morphismes de monoïdes $u:\mathrm{G}\to\mathrm{G}'$ sur l'ensemble des foncteurs
\[ \mathrm{F}:\ \mathrm{Rep}(\mathrm{G}')\to\mathrm{Rep}(\mathrm{G}) \]
tels que $\omega^{\mathrm{G}}\circ\mathrm{F}=\omega^{\mathrm{G}'}$.

Théorème 1.5. --- Soient S un schéma affine, G un S-groupe affine et plat,
et soient $\omega$, $\omega'$, des foncteurs fibre sur Rep(G). On a alors :

(a) $\mathit{Isom}^{\otimes}(\omega,\omega')=\mathit{Hom}^{\otimes}(\omega,\omega')$ ;

(b) $\omega$ et $\omega'$ sont localement isomorphes pour la topologie fpqc, i.e., il
existe un recouvrement $\{\mathrm{S}_i\to\mathrm{S}\}$ de S au sens de la topologie fpqc, et des
$\otimes$-isomorphismes $\omega_{\mathrm{S}_i}\simeq\omega'_{\mathrm{S}_i}$.

Si on note Tors(G) la catégorie des torseurs à droite sous G [($^5$), chap. II],
on peut déduire du théorème qui précède et de 1.3 :

Corollaire 1.6. --- La correspondance $\omega\mapsto\mathrm{Isom}^{\otimes}(\omega^{\mathrm{G}},\omega)$ induit une
équivalence de catégories
\[ (1.6.1)\qquad \mathrm{Fib}(\mathrm{G})\to\mathrm{Tors}(\mathrm{G}), \]
et en particulier une bijection
\[ (1.6.2)\qquad \mathrm{fib}(\mathrm{G})\to\mathrm{H}^1(\mathrm{S}_{\mathrm{fpqc}},\mathrm{G}), \]
où $\mathrm{S}_{\mathrm{fpqc}}$ désigne le site des S-schémas munis de la topologie fpqc.

Remarque 1.7. --- Si le S-groupe affine et plat G est de présentation
finie (resp. lisse), on peut dans l'assertion (b) de 1.5 remplacer la topo-%
logie fpqc par la topologie fppf (resp. étale).

2. Conditions de finitude.

2.1. Il s'agit ici d'exhiber des sous-$\otimes$-catégories de Rep(G), qui soient
de nature « finie », et qui donnent des résultats semblables à ceux présentés
au paragraphe précédent. On va se limiter au cas où S est le spectre d'un
corps $k$, où les résultats sont plus sympathiques.

Soit G un $k$-monoïde affine. On note $\mathrm{Rep}_0(\mathrm{G})$ la sous-$\otimes$-catégorie pleine
de Rep(G) des G-$k$-modules de rang fini sur $k$, $\omega_0^{\mathrm{G}}:\mathrm{Rep}(\mathrm{G})\to\mathrm{Modf}(k)$
le foncteur « oubli de l'action de G ». Si T est un $k$-schéma, et on note
Loclib(T) la $\otimes$-catégorie ACU $k$-linéaire des $\mathcal{O}_{\mathrm{T}}$-Modules localement libres
de rang fini, on appelle foncteur fibre sur $\mathrm{Rep}_0(\mathrm{G})$ à valeurs dans T un
$\otimes$-foncteur ACU $k$-linéaire $\omega:\mathrm{Rep}_0(\mathrm{G})\to\mathrm{Loclib}(\mathrm{T})$ qui soit fidèle et
exact. On obtient encore une catégorie $\mathrm{Fib}_0(\mathrm{G},\mathrm{T})$.

\mspage{\hand{70}}
